%% frontmatter/dependency-graph.tex
%% Chapter dependency graph — rebuilt to pass the visual audit per STYLE.md §10.1.
%%
%% Layout choice (deliberate): one row per Part, with arrows only WITHIN
%% each row. Reading top-to-bottom implies progression between Parts; the
%% Part labels on the left make this explicit. This eliminates the long
%% cross-row arrows that broke the previous version.

\chapter*{How the chapters depend on each other}
\addcontentsline{toc}{chapter}{Chapter dependency graph}

The chart below is a reading roadmap. Each row is one Part; within each row, an arrow ``\textbf{Ch.\,A} $\to$ \textbf{Ch.\,B}'' means \emph{read Ch.\,A first}. Read top-to-bottom to progress through the book.

If you only have time for the spine, follow the cliff-notes track:
\begin{quote}
  Ch.\,1 $\to$ Ch.\,4 $\to$ Ch.\,7 $\to$ Ch.\,10 $\to$ Ch.\,12 $\to$ Ch.\,17 $\to$ Ch.\,21.
\end{quote}

\bigskip

\begin{fullwidth}
\centering
\begin{tikzpicture}[
  >={Stealth[length=4pt]},
  every node/.style={font=\sffamily},
  chap/.style={
    rectangle, rounded corners=3pt,
    draw=black!60, line width=0.4pt,
    fill=mathemagicsDefBg,
    inner sep=4pt,
    minimum width=2.0cm, minimum height=0.95cm,
    align=center,
    font=\footnotesize\sffamily
  },
  spine/.style={
    rectangle, rounded corners=3pt,
    draw=mathemagicsDef, line width=0.7pt,
    fill=mathemagicsDefBg,
    inner sep=4pt,
    minimum width=2.0cm, minimum height=0.95cm,
    align=center,
    font=\footnotesize\sffamily\bfseries
  },
  needs/.style={->, semithick, draw=black!75, line width=0.5pt,
                shorten >=2pt, shorten <=2pt},
  partlbl/.style={font=\bfseries\small\sffamily, align=right,
                  anchor=east, color=black!65}
]

% ----------------------------------------------------------------
% Row 1: Part I — Foundations of Probability
% ----------------------------------------------------------------
\node[partlbl] at (-0.7, 0)             {Part~I};
\node[spine]   at ( 0.5, 0)        (c1) {Ch.\,1\\RVs};
\node[chap]    at ( 3.0, 0)        (c2) {Ch.\,2\\MGF/CF};
\node[chap]    at ( 5.5, 0)        (c3) {Ch.\,3\\Conv.};
\node[spine]   at ( 8.0, 0)        (c4) {Ch.\,4\\CLT};

\draw[needs] (c1) -- (c2);
\draw[needs] (c2) -- (c3);
\draw[needs] (c3) -- (c4);

% ----------------------------------------------------------------
% Row 2: Part II — From Numbers to Vectors and Matrices
% ----------------------------------------------------------------
\node[partlbl] at (-0.7, -1.6)            {Part~II};
\node[chap]    at ( 0.5, -1.6)      (c5) {Ch.\,5\\Vectors};
\node[chap]    at ( 3.0, -1.6)      (c6) {Ch.\,6\\Linalg};
\node[spine]   at ( 5.5, -1.6)      (c7) {Ch.\,7\\PSD};

\draw[needs] (c5) -- (c6);
\draw[needs] (c6) -- (c7);

% ----------------------------------------------------------------
% Row 3: Part III — Sample Covariance Matrices
% ----------------------------------------------------------------
\node[partlbl] at (-0.7, -3.2)             {Part~III};
\node[chap]    at ( 0.5, -3.2)       (c8) {Ch.\,8\\SCM};
\node[chap]    at ( 3.0, -3.2)       (c9) {Ch.\,9\\Eig.\,conv.};
\node[spine]   at ( 5.5, -3.2)      (c10) {Ch.\,10\\Wishart};

\draw[needs] (c8) -- (c9);
\draw[needs] (c9) -- (c10);

% ----------------------------------------------------------------
% Row 4: Part IV — Eigenvalue Fluctuations in Fixed Dimension
% ----------------------------------------------------------------
\node[partlbl] at (-0.7, -4.8)             {Part~IV};
\node[chap]    at ( 0.5, -4.8)      (c11) {Ch.\,11\\Joint CLT};
\node[spine]   at ( 3.0, -4.8)      (c12) {Ch.\,12\\Two regimes};
\node[chap]    at ( 5.5, -4.8)      (c13) {Ch.\,13\\Eigvecs};
\node[chap]    at ( 8.0, -4.8)      (c14) {Ch.\,14\\GOE};

\draw[needs] (c11) -- (c12);
\draw[needs] (c12) -- (c13);
\draw[needs] (c13) -- (c14);

% ----------------------------------------------------------------
% Row 5: Part V — High-Dimensional Phenomena
% ----------------------------------------------------------------
\node[partlbl] at (-0.7, -6.4)             {Part~V};
\node[chap]    at ( 0.5, -6.4)      (c15) {Ch.\,15\\High-dim};
\node[chap]    at ( 3.0, -6.4)      (c16) {Ch.\,16\\MP};
\node[spine]   at ( 5.5, -6.4)      (c17) {Ch.\,17\\TW};
\node[chap]    at ( 8.0, -6.4)      (c18) {Ch.\,18\\Q-Q};

\draw[needs] (c15) -- (c16);
\draw[needs] (c16) -- (c17);
\draw[needs] (c17) -- (c18);

% ----------------------------------------------------------------
% Row 6: Part VI — Spikes and Spectral Clustering
% ----------------------------------------------------------------
\node[partlbl] at (-0.7, -8.0)             {Part~VI};
\node[chap]    at ( 0.5, -8.0)      (c19) {Ch.\,19\\Spikes};
\node[chap]    at ( 3.0, -8.0)      (c20) {Ch.\,20\\BBP};
\node[spine]   at ( 5.5, -8.0)      (c21) {Ch.\,21\\Spec.\,Clust.};

\draw[needs] (c19) -- (c20);
\draw[needs] (c20) -- (c21);

% ----------------------------------------------------------------
% Legend
% ----------------------------------------------------------------
\node[font=\footnotesize\itshape\sffamily, anchor=west, color=black!65]
  at (-0.7, -9.4)
  {Bordered nodes are on the cliff-notes spine.};

\end{tikzpicture}
\end{fullwidth}

\bigskip

\noindent\textit{The graph is intentionally lossy --- it shows the within-Part reading order, not every cross-reference. When in doubt, the table of contents always works.}
